\chapter{变换光学与全息照相}
\problem{5-1}{设薄透镜由折射率为 $n_0$ 的材料做成，物方和像方的折射率分别是 $n$ 和 $n'$，导出其相位变换函数（用透镜的焦距表示出来）。}
\answerarea
\problem{5-2}{用薄透镜的相位变换函数导出傍轴条件下的横向放大率公式。}
\answerarea
\clearpage
\problem{5-3}{用劈形棱镜的相位变换函数（5.27）式导出傍轴光束斜入射时产生的偏向角 $\delta$。
\begin{equation*}
\widetilde T_p(x,y)=\exp\!\left[-\mathrm j k(n-1)\alpha x\right].\tag{5.27}
\end{equation*}
其中 $\mathrm j^2=-1$，$k$ 为入射光波数，$n$ 为棱镜折射率，$\alpha$ 为小劈角。}
\answerarea
\problem{5-4}{在一相干成像系统中，镜头（作为入射光瞳）的相对孔径为 $1/5$，求此系统的截止频率（$\mathrm{mm}^{-1}$）。设物平面在前焦面附近，照明波长为 $0.5\,\mu\mathrm m$。}
\answerarea
\clearpage
\problem{5-5}{利用阿贝成像原理导出在相干照明条件下显微镜的最小分辨距离公式，并同非相干照明的最小分辨距离公式比较。}
\answerarea
\problem{5-6}{采用远场装置（图5-20b）接收单缝的夫琅禾费衍射场，设单缝宽约为 $100\,\mu\mathrm m$，入射光波长 $632.8\,\mathrm{nm}$。
\begin{enumerate}\item 接收屏幕至少应放在多远？\item 在接收屏幕的多大范围内才算是夫琅禾费衍射场？\item 0级半角宽度为多少？\item 在接收屏幕上0级的线宽度有多少？\end{enumerate}
\bookfig[0.92]{book_5_20b}{图5-20(b)\quad 远场接收}}
\answerarea
\clearpage
\problem{5-7}{采用像面接收装置（图5-20d或e）接收单缝的夫琅禾费衍射场，设单缝宽度约为 $1\,\mathrm{mm}$，入射光波长 $488.0\,\mathrm{nm}$，物距 $40\,\mathrm{cm}$，像距 $80\,\mathrm{cm}$。
\begin{enumerate}\item 如果单缝置于透镜后方，要求在像面 $1\,\mathrm{cm}$ 范围内准确地接收到夫琅禾费衍射场，单缝离像面至少多远？\item 如果单缝紧贴透镜后面，求0级半角宽度和接收屏幕上0级的线宽度；\item 如果单缝离透镜 $40\,\mathrm{cm}$ 远，求0级半角宽度及它在幕上的线宽度；\item 如果单缝置于透镜前方，紧贴在其左侧，情况如何？\end{enumerate}
\bookfig[0.88]{book_5_20d}{图5-20(d)\quad 像面接收（一）}
\bookfig[0.88]{book_5_20e}{图5-20(e)\quad 像面接收（二）}}
\answerarea
\clearpage
\problem{5-8}{对图5-20e所示装置，推导出傍轴条件下它的衍射场表达式，并论证它符合夫琅禾费衍射场的标准形式。【提示：利用场点 $P'$ 的共轭点 $P'_c$（见图5-20e）。】
\bookfig[0.88]{book_5_20e}{图5-20(e)\quad 像面接收（二）}}
\answerarea
\problemfig[0.17]{5-9}{在透镜的前焦面上有一系列同相位的相干光源等距排列在 $x$ 轴上，形成一维点阵（见本题图），用傅里叶变换法求后焦面上的夫琅禾费衍射场。}{5_9}
\answerarea
\clearpage
\problem{5-10}{设透镜直径 $D=5\,\mathrm{cm}$，焦距 $F=60\,\mathrm{cm}$，图像（衍射屏）线度 $l=2\,\mathrm{cm}$，入射光波长 $\lambda=0.6\,\mu\mathrm m$。
\begin{enumerate}
\item 分别算出与后焦面上下列各点 $(x',y')$ 相对应的空间频率 $(f_x,f_y)$ 的具体数值：
\[
\begin{gathered}
(0,0),\quad(0,1),\quad(1,0),\quad
\left(\frac{\sqrt2}{2},\frac{\sqrt2}{2}\right),\\[4pt]
\underbrace{\left(-\frac{\sqrt2}{2}\right)}_{\text{原文坐标缺项}},\quad
(0.5,2),\quad(3,-5),\quad(-10,-15).
\end{gathered}
\]
坐标单位皆为 $\mathrm{mm}$，空间频率单位皆为 $\mathrm{mm}^{-1}$。
\item 计算系统的截止频率。
\end{enumerate}
\editorial{按补充的清晰截图核对：第五项原文仅印为 $(-\sqrt2/2)$，缺少一个坐标分量，照录并标注；其后一项明确为 $(0.5,2)$。}}
\answerarea
\clearpage
\problem{5-11}{如本题图a，参考光束 $R$ 和物光束 $O$ 均为平行光，对称地斜射于记录介质 $\Sigma$ 上，即 $\theta_R=-\theta_O$，二者间的夹角 $\theta=2\theta_O$。
\begin{enumerate}\item 试说明全息图上干涉条纹的形状；\item 试分别写出物光波和参考光波在 $\Sigma$ 平面上的相位分布 $\varphi_O(y)$ 和 $\varphi_R(y)$；\item 试证明全息图上干涉条纹的间距公式为 $d=\lambda/[2\sin(\theta/2)]$；\item 试计算，当夹角 $\theta=1^\circ$ 时间距 $d$ 为多少？当夹角 $\theta=60^\circ$ 时间距 $d$ 为多少？（采用 He--Ne 激光记录，$\lambda=632.8\,\mathrm{nm}$。）\item 某感光胶片厂生产一种可用于全息照相的记录干版，其性能如下：感光层厚度为 $8\,\mu\mathrm m$，分辨率为3000条/mm以上。利用题（4）所得数据，试说明：当夹角 $\theta=60^\circ$ 时，用该记录干版是否构成一张体全息图？当夹角 $\theta=60^\circ$ 时，该记录干版的分辨率是否匹配？\item 如本题图b，采用与参考光束 $R$ 同样波长同样倾角的照明光束 $R'$ 照射这张全息图，试分析0级、$+1$级、$-1$级三个衍射波都出现在什么方向上，并在图上画出。\end{enumerate}\fig[0.85]{5_11}
\editorial{原材料将此配图标为“习题5-12”；5-11与5-12共用此图。}}
\answerarea
\clearpage
\problem{5-12}{若在上题中改为用正入射的平面波再现，$\pm1$ 级衍射波各发生什么变化？}
\answerarea
\problem{5-13}{如本题图a，用正入射的平面参考光波记录轴外物点 $O$ 发出的球面波。
\begin{enumerate}\item 如本题图b所示，用轴上的点源 $R'$ 发出的球面波来重建波前，求 $\pm1$ 级两像点的位置；\item 用与记录全息图时不同波长的正入射平面波照明，求 $\pm1$ 级两像点的位置。\end{enumerate}\fig[0.73]{5_13}}
\answerarea
\clearpage
\problem{5-14}{（1）求图5-32a所示装置制备的全息图中各级干涉条纹的半径，并证明它们与一张菲涅耳波带片相符；（2）验证用图5-32c方式再现的两个像点 $O$、$O'$ 确是菲涅耳波带片的一对焦点。
\bookfig[0.96]{book_5_32}{图5-32\quad 菲涅耳波带片与全息图}
\par{\small 图（a）用轴上点源 $O$ 的球面波和正入射平面参考波 $R$ 记录全息图；图（b）为同心圆干涉条纹示意；图（c）用正入射平面波 $R'$ 再现，形成像点 $O$、$O'$。}}
\answerarea
\clearpage
